% généré par scripts/06_textes.py
\dosmeta{Année de classement : 348 (directe : nom consulaire lu). Références : CIL VI 33865 = CLE 659 = ICUR I 887 = ILCV 748 (add). Rattachement selon l’EDB : \emph{Incertae originis} (le volume I des ICVR réunit les inscriptions dont la provenance est perdue). Support : Texte métrique. Signes décrits par l’EDB : chrisme et chrisme radié. Reproductions signalées par l’EDB : dessin, ICVR I (de Rossi) (droits réservés ; non reproduites). Fiche EDB : Filippo Piazzolla, première version 2014. Notice du catalogue p.~\pageref{not:22}.}
\txtbloc{Texte transmis (EDCS, état 2022)}{\lnum{1}Felix vita viri felicior exitus ipse\par
\lnum{2}Caiani semper crescit per saecula nomen\par
\lnum{3}nescit fama mori sed semper vi⸢v⸣it ubique\par
\lnum{4}advenit hospes Romanus princeps in urbem\par
\lnum{5}qui fuit hic primum iuris consultor amicus\par
\lnum{6}quiescit in pace\{m\} depositus diem quar\par
\lnum{7}tum Nonas Augu\par
\lnum{8}stas Flavio Filippo\par
\lnum{9}et Flavio Sallea con\par
\lnum{10}sulibus pater Sabba\par
\lnum{11}tius fecit}
\txtbloc{Texte de l’EDB (EDB29417)}{\lnum{1}felix vita viri felicior exitus ipse\par
\lnum{2}Caiani semper crescit per saecula nomen\par
\lnum{3}nescit fama mori sed semper vibit ubique\par
\lnum{4}advenit hospes romanus princeps in urbem\par
\lnum{5}qui fuit hic primum iuris consultor amicus\par
\lnum{6}quiescit in pacem depositus diem quar/tum nonas augu/stas Flavio Filippo\par
\lnum{7}et Flavio Sallea con/sulibus pater Sabba/tius fecit ((avis)) ((monogramma Christi radiatum))}
